\section{第一讲：Springer 纤维}

Springer 纤维是几何表示论中经典而基本的对象。Springer \cite{Springer} 首先发现，它们的上同调群实现 Weyl 群的表示，这一现象称为 Springer 对应。作为奇异的代数簇，Springer 纤维本身即是有趣的几何对象。经由特征层，它们还与有限 Lie 型群的表示论相联系。

\subsection{设定}\label{ss:notation} 本节中 $k$ 为代数闭域。设 $G$ 是 $k$ 上连通约化群，其伴随群单（于是伴随群 $G^{\ad}$ 由七族 Dynkin 图中的一族确定）。设 $\chark$ 相对 $G$ 足够大。记 $r$ 为 $G$ 的秩。

有时方便固定 $\frg$ 的一个 Cartan 子代数 $\frt$，等价地固定一个极大环面 $T\subset G$。一旦如此，就可以谈论 $T$ 在 $\frg$ 上作用的根。

设 $\sB$ 为 $G$ 的旗簇：这是参数化 $G$ 的 Borel 子群的 $G$ 齐性射影簇。取定 Borel 子群 $B\subset G$，可把 $\sB$ 与 $G/B$ 等同。

记 $\frg=\Lie G$。对 $X\in\frg$，记 $C_{G}(X)$ 为 $X$ 在 $G$ 中的中心化子，即 $G$ 在 $\frg$ 上伴随作用下 $X$ 的稳定化子。

设 $\sN\subset\frg$ 为幂零元构成的子簇。这是一个锥：它在 $\Gm$ 按缩放作用在 $\frg$ 上时稳定。已知在伴随作用下，$\sN$ 上只有有限多个 $G$ 轨道。

\subsection{Springer 纤维}

\subsubsection{Springer 消解} 余切丛 $\tN:=T^{*}\sB$ 参数化这样的数对 $(e, B)$：$e\in\sN$，$B$ 是 $G$ 的 Borel 子群，且 $e\in \frn$，其中 $\frn$ 是 $\Lie B$ 的幂零根。{\em Springer 消解}是遗忘映射
\begin{equation*}
\pi:\tN\to\sN
\end{equation*}
它把 $(e, B)$ 送到 $e$。此映射是射影的。

对 $e\in\sN$，纤维 $\sB_{e}:=\pi^{-1}(e)$ 称为 $e$ 的 {\em Springer 纤维}。按定义，$\sB_{e}$ 是 $\sB$ 中由那些满足「$e$ 落在 $\Lie B$ 的幂零根中」的 Borel 子群 $B\subset G$ 构成的闭子概形。

\subsubsection{Grothendieck 修改}\label{sss:Gr alt}
考虑由数对 $(X,B)$ 构成的簇，其中 $X\in\frg$，$B\in\sB$，且 $X\in\Lie B$。遗忘映射 $(X,B)\mapsto X$
\begin{equation*}
\pi_{\frg}: \tg\to \frg
\end{equation*}
称为 {\em Grothendieck 修改} \footnote{术语「修改」（alteration）指源在 $k$ 上光滑的真、一般有限的映射。}，又称 {\em Grothendieck 同时消解}\zy{两名为同一映射 $\pi_{\frg}:\tg\to\frg$ 的两种通行叫法，不是两个对象：以真、一般有限的覆盖视之，称「修改」；以幂零锥上的消解视之，称「同时消解」。来源：原文 \S1.2 该段文脉（\texttt{ZhiweiYunPCMI.tex} 第 935--951 行，含作者脚注）。}。

设 $\frc=\frg\sslash G\cong \frt\sslash W$ 为 $\frg$ 在 $G$ 伴随作用下的范畴商。于是有交换图
\begin{equation}\label{tgg}
\xymatrix{\tg\ar[d]^{\pi_{\frg}}\ar[r]^{\wt\chi} & \frt\ar[d]\\
\frg \ar[r]^{\chi} & \frc}
\end{equation}
其中 $\chi:\frg\to \frc$ 是自然商映射，$\wt\chi:\tg\to \frt$ 把 $(X,B)\in\tg$ 送到 $X$ 在 $\frb/\frn$ 中的像（这里 $\frb=\Lie B$，幂零根为 $\frn$）；选取含 $\frt$ 的 Borel 子群后，$\frb/\frn$ 可与 $\frt$ 典范等同。把左列限制到 $\frg$ 的正则元上时，图 \eqref{tgg} 是笛卡儿方块 \footnote{若 $X\in\frg$ 在 $G$ 中的中心化子维数为 $r$（$G$ 的秩），则称 $X$ 为{\em 正则}元。}。特别地，若把图 \eqref{tgg} 限制到正则半单轨迹 $\frc^{\rs}\subset \frc$ \footnote{若对任意根 $\alpha$ 都有 $\alpha(X)\neq0$，则称 $X\in\frt$ 为{\em 正则半单}元。用 $\frt^{\rs}\subset\frt$ 记 $\frt$ 中正则半单元构成的开子概形。由于 $\frt^{\rs}$ 在 $W$ 下稳定，它是某个开子集 $\frc^{\rs}\subset\frc$ 的原像。开子集 $\frg^{\rs}=\chi^{-1}(\frc^{\rs})$ 按定义就是 $\frg$ 的{\em 正则半单}轨迹。}，则 $\pi_{\frg}|_{\tg^{\rs}}:\tg^{\rs}\to \frg^{\rs}$ 是 $W$-挠子；也就是说，$W$ 作用在 $\pi_{\frg}^{-1}(\frg^{\rs})$ 上，保持到 $\frg^{\rs}$ 的投影，并使 $\pi_{\frg}$ 的纤维成为 $W$ 的主齐性空间。在正则轨迹 $\frg^{\reg}\subset \frg$ 上方，$\pi_{\frg}$ 出现分支，但仍有限。

对 $X\in \frg$，把纤维 $\pi_{\frg}^{-1}(X)$ 简记为 $\tB_{X}$。把 $\pi_{\frg}$ 限制到 $\sN$ 上，幂零元 $e\in\sN$ 对应的纤维 $\tB_{e}$ 是 $\sB$ 中由满足 $e\in\Lie B$ 的 $B$ 构成的闭子簇。显然 $\sB_{e}$ 是 $\tB_{e}$ 的子概形，且两者有相同的约化结构。然而作为概形，$\tB_{e}$ 与 $\sB_{e}$ 一般不同。参见下文 \S\ref{sss:regular Spr} 与习题 \ref{exI:nonred Spr}。
